% GENERATED FILE. Edit canonical lessons, then run npm run generate:books.
% canonical-source-hash: fnv1a64:8034ca24bf347cb5
% canonical-lessons: GU-C526-sur, GU-C526-tal, GU-C526-gayak, GU-C526-rag, GU-C526-chedvu

\chapter{Notes and Rhythm}
\label{ch:gu-notes-and-rhythm}
\hlchaptermodality{526}

\begin{chapteropening}
\textbf{By the end of this chapter:} \emph{I can talk about notes, rhythm, ragas and singers.}

Complete the last lesson of chapter 526: I can talk about notes, rhythm, ragas and singers.
\end{chapteropening}

\section[\texorpdfstring{\gu{સૂર} (sūr)}{સૂર (sūr)}]{\gu{સૂર} (sūr) — a musical note}

\label{lesson:GU-C526-sur}



\begin{quote}
Before the new one: say the Gujarati for to frame (a picture), then the Gujarati for an exhibition.
\end{quote}

\subsection*{You'll want to know: \gu{સૂર}}
\textbf{\gu{સૂર}} — \emph{sūr} — \textquotedblleft{}a musical note\textquotedblright{}.

\begin{quote}
\gu{મીરાનો} \gu{સૂર} \gu{ઘણો} \gu{મીઠો} \gu{છે}.
\end{quote}

\begin{quote}
\emph{Mīrāno sūr ghaṇo mīṭho che.}
\end{quote}

\begin{quote}
Mira sings very sweetly.
\end{quote}

A masculine noun.

\begin{cousinweb}
\textbf{\gu{સૂર}} (\emph{sūr}) is filed by CDIAL under Sanskrit \emph{sura}, and CDIAL lists the Gujarati word as \textquotedblleft{}a modulation of the voice, a tune\textquotedblright{} (CDIAL 13498).
\end{cousinweb}

\begin{grammarlens}[title={What you've built}]
One more word for describing what you lived and felt.
\end{grammarlens}

\subsection*{Your turn}
\begin{itemize}
\raggedright
  \item \emph{Say it:} \emph{sūr}
  \item \emph{Say it:} \emph{sūr}, once more
  \item \emph{Say it:} \emph{Mīrāno sūr ghaṇo mīṭho che.} — Mira sings very sweetly.
  \item \emph{Recall:} say the Gujarati for to frame (a picture), then the Gujarati for an exhibition, then say \emph{sūr} again
\end{itemize}

\begin{tcolorbox}[breakable,colback=teal!4,colframe=teal!35!black,title={Before you move on}]
What does \gu{સૂર} mean? (\textquotedblleft{}A musical note\textquotedblright{}.) Say it once more.
\end{tcolorbox}
\section[\texorpdfstring{\gu{તાલ} (tāl)}{તાલ (tāl)}]{\gu{તાલ} (tāl) — a rhythm (a beat cycle)}

\label{lesson:GU-C526-tal}



\begin{quote}
Before the new one: say the Gujarati for an exhibition, then the Gujarati for a musical note.
\end{quote}

\subsection*{You'll want to know: \gu{તાલ}}
\textbf{\gu{તાલ}} — \emph{tāl} — \textquotedblleft{}a rhythm (a beat cycle)\textquotedblright{}.

\begin{quote}
\gu{તબલાંનો} \gu{તાલ} \gu{ઘણો} \gu{ઝડપી} \gu{હતો}.
\end{quote}

\begin{quote}
\emph{tablā̃no tāl ghaṇo jhaḍpī hato.}
\end{quote}

\begin{quote}
The rhythm of the tabla was very fast.
\end{quote}

A masculine noun.

\begin{cousinweb}
\textbf{\gu{તાલ}} (\emph{tāl}) is a tatsama, taken from Sanskrit \emph{tāla}, which Monier-Williams gives for the palm tree and also for a musical beat (Monier-Williams).
\end{cousinweb}

\begin{grammarlens}[title={What you've built}]
One more word for describing what you lived and felt.
\end{grammarlens}

\subsection*{Your turn}
\begin{itemize}
\raggedright
  \item \emph{Say it:} \emph{tāl}
  \item \emph{Say it:} \emph{tāl}, once more
  \item \emph{Say it:} \emph{tablā̃no tāl ghaṇo jhaḍpī hato.} — The rhythm of the tabla was very fast.
  \item \emph{Recall:} say the Gujarati for an exhibition, then the Gujarati for a musical note, then say \emph{tāl} again
\end{itemize}

\begin{tcolorbox}[breakable,colback=teal!4,colframe=teal!35!black,title={Before you move on}]
What does \gu{તાલ} mean? (\textquotedblleft{}A rhythm (a beat cycle)\textquotedblright{}.) Say it once more.
\end{tcolorbox}
\section[\texorpdfstring{\gu{ગાયક} (gāyak)}{ગાયક (gāyak)}]{\gu{ગાયક} (gāyak) — a singer}

\label{lesson:GU-C526-gayak}



\begin{quote}
Before the new one: say the Gujarati for a musical note, then the Gujarati for a rhythm.
\end{quote}

\subsection*{You'll want to know: \gu{ગાયક}}
\textbf{\gu{ગાયક}} — \emph{gāyak} — \textquotedblleft{}a singer\textquotedblright{}.

\begin{quote}
\gu{આ} \gu{ગાયકનાં} \gu{ગીતો} \gu{મને} \gu{ઘણાં} \gu{ગમે} \gu{છે}.
\end{quote}

\begin{quote}
\emph{ā gāyaknā̃ gīto mane ghaṇā̃ game che.}
\end{quote}

\begin{quote}
I like this singer's songs very much.
\end{quote}

A masculine or feminine noun, by the person meant.

\begin{cousinweb}
\textbf{\gu{ગાયક}} (\emph{gāyak}) is a tatsama, taken from Sanskrit \emph{gāyaka}, \textquotedblleft{}a singer\textquotedblright{} (Monier-Williams), from the root of \textbf{\gu{ગાવું}}, to sing.
\end{cousinweb}

\begin{grammarlens}[title={What you've built}]
One more word for describing what you lived and felt.
\end{grammarlens}

\subsection*{Your turn}
\begin{itemize}
\raggedright
  \item \emph{Say it:} \emph{gāyak}
  \item \emph{Say it:} \emph{gāyak}, once more
  \item \emph{Say it:} \emph{ā gāyaknā̃ gīto mane ghaṇā̃ game che.} — I like this singer's songs very much.
  \item \emph{Recall:} say the Gujarati for a musical note, then the Gujarati for a rhythm, then say \emph{gāyak} again
\end{itemize}

\begin{tcolorbox}[breakable,colback=teal!4,colframe=teal!35!black,title={Before you move on}]
What does \gu{ગાયક} mean? (\textquotedblleft{}A singer\textquotedblright{}.) Say it once more.
\end{tcolorbox}
\section[\texorpdfstring{\gu{રાગ} (rāg)}{રાગ (rāg)}]{\gu{રાગ} (rāg) — a raga (a melodic mode)}

\label{lesson:GU-C526-rag}



\begin{quote}
Before the new one: say the Gujarati for a rhythm, then the Gujarati for a singer.
\end{quote}

\subsection*{You'll want to know: \gu{રાગ}}
\textbf{\gu{રાગ}} — \emph{rāg} — \textquotedblleft{}a raga (a melodic mode)\textquotedblright{}.

\begin{quote}
\gu{ગાયકે} \gu{સાંજે} \gu{એક} \gu{જૂનો} \gu{રાગ} \gu{ગાયો}.
\end{quote}

\begin{quote}
\emph{gāyke sā̃je ek jūno rāg gāyo.}
\end{quote}

\begin{quote}
In the evening the singer sang an old raga.
\end{quote}

A masculine noun.

\begin{cousinweb}
\textbf{\gu{રાગ}} (\emph{rāg}) is a tatsama, taken from Sanskrit \emph{rāga}, \textquotedblleft{}colouring\textquotedblright{}, which Monier-Williams also gives for a musical mode (Monier-Williams).
\end{cousinweb}

\begin{grammarlens}[title={What you've built}]
One more word for describing what you lived and felt.
\end{grammarlens}

\subsection*{Your turn}
\begin{itemize}
\raggedright
  \item \emph{Say it:} \emph{rāg}
  \item \emph{Say it:} \emph{rāg}, once more
  \item \emph{Say it:} \emph{gāyke sā̃je ek jūno rāg gāyo.} — In the evening the singer sang an old raga.
  \item \emph{Recall:} say the Gujarati for a rhythm, then the Gujarati for a singer, then say \emph{rāg} again
\end{itemize}

\begin{tcolorbox}[breakable,colback=teal!4,colframe=teal!35!black,title={Before you move on}]
What does \gu{રાગ} mean? (\textquotedblleft{}A raga (a melodic mode)\textquotedblright{}.) Say it once more.
\end{tcolorbox}
\section[\texorpdfstring{\gu{છેડવું} (cheḍvũ)}{છેડવું (cheḍvũ)}]{\gu{છેડવું} (cheḍvũ) — to strike up (a tune)}

\label{lesson:GU-C526-chedvu}



\begin{quote}
Before the new one: say the Gujarati for a singer, then the Gujarati for a raga.
\end{quote}

\subsection*{You'll want to know: \gu{છેડવું}}
\textbf{\gu{છેડવું}} — \emph{cheḍvũ} — \textquotedblleft{}to strike up (a tune)\textquotedblright{}.

\begin{quote}
\gu{મીરાએ} \gu{વાંસળી} \gu{પર} \gu{એક} \gu{મીઠો} \gu{સૂર} \gu{છેડ્યો}.
\end{quote}

\begin{quote}
\emph{Mīrāe vā̃sḷī par ek mīṭho sūr cheḍyo.}
\end{quote}

\begin{quote}
Mira struck up a sweet tune on the flute.
\end{quote}

Past, with \gu{તેણે} (he or she …): \gu{તેણે} \gu{સૂર} \gu{છેડ્યો}, \emph{teṇe sūr cheḍyo}.

\begin{cousinweb}
\textbf{\gu{છેડવું}} (\emph{cheḍvũ}) is filed by CDIAL under \emph{*ciḍ}, \textquotedblleft{}to be angry\textquotedblright{}, and CDIAL lists the Gujarati verb as \textquotedblleft{}to irritate\textquotedblright{} (CDIAL 4794); Gujarati also uses it for touching off a tune on an instrument.
\end{cousinweb}

\begin{grammarlens}[title={What you've built}]
One more word for describing what you lived and felt.
\end{grammarlens}

\subsection*{Your turn}
\begin{itemize}
\raggedright
  \item \emph{Say it:} \emph{cheḍvũ}
  \item \emph{Say it:} \emph{cheḍvũ}, once more
  \item \emph{Say it:} \emph{Mīrāe vā̃sḷī par ek mīṭho sūr cheḍyo.} — Mira struck up a sweet tune on the flute.
  \item \emph{Recall:} say the Gujarati for a singer, then the Gujarati for a raga, then say \emph{cheḍvũ} again
\end{itemize}

\begin{tcolorbox}[breakable,colback=teal!4,colframe=teal!35!black,title={Before you move on}]
What does \gu{છેડવું} mean? (\textquotedblleft{}To strike up (a tune)\textquotedblright{}.) Say it once more.
\end{tcolorbox}
